\section{HP-MR transient: +0.5\,\$ external reactivity insertion}
\label{sec:transient}

\paragraph{Reference.} The reference is $P(t)$ from \garnet{} on the GH200.
Fixed-step BDF2 at $\Delta t$ = 0.025 and 0.0125~s is extrapolated in time
(order 2), then in space over $(r, n_z)$ = (4, 50), (6, 80) and (8, 100). The
two corrections are $1.5 \times 10^{-4}$ and $6.8 \times 10^{-5}$ relative,
so the reference is good to about $10^{-4}$.

The ($r=12$, $n_z=160$) mesh would have tightened it, but it does not fit in
the GH200's 96~GB: the transient's precursor fields and BDF history raised
the allocation past 100~GB, where the eigen-solve on the same mesh needs
74~GB. \ffn's own runs converge onto this reference independently: its
Semi-Implicit-Exponential error falls from $1.1 \times 10^{-3}$ to
$4.5 \times 10^{-4}$ as $\Delta t$ halves to 0.003125~s. $P(2\,\mathrm{s}) =
2.5206$.

\paragraph{Time to accuracy.} See Table~\ref{tab:tta} (last rows) and
Figure~\ref{fig:transient}.
\begin{itemize}
\item \textbf{$10^{-2}$:} the two codes are even on one core (\ffn{} 107~s,
  \garnet{} 100~s). The GH200 takes 12~s.
\item \textbf{$10^{-3}$:} \garnet{} on one core needs 127~s ($r=2$, BDF2,
  $\Delta t = 0.1$~s) and \ffn{} 1262~s (FE2 $r=0$, $\Delta t = 0.003125$~s),
  a factor of 10. On the GH200 it is 17~s, 76 times faster than \ffn.
\end{itemize}

The \garnet{} transients were not rerun with the Wielandt shift
(Section~\ref{sec:steady}). Their initial eigen-solve is the power iteration,
the same solve as the steady $r=2$ run: 18.2~s on one core and 3.0~s on the
GH200, against 8.3 and 1.0~s with the shift. At $10^{-3}$ the shift would
therefore save about 10~s of the 127~s on one core and 2~s of the 17~s on the
GH200; the comparison with \ffn{} does not change.

The difference is the time integration, not the space discretization. \ffn's
two working schemes are first order. Its default Implicit-Exponential scheme
does not hold the unperturbed state on this input
(Section~\ref{sec:problems}). \garnet's BDF2 is second order. \ffn's spatial
error is small at every mesh: FE2 $r=0$ and FE3 $r=0$ differ by less than
$10^{-3}$. Its error is therefore almost all time error, and a $10^{-3}$
target costs it 640 steps.

\paragraph{\garnet{} step control.} Down to about $10^{-3}$ the front is
fixed-step BDF2 with the fixed-point step. Below that it is adaptive BDF with
automatic order up to BDF4 and an error per step of $10^{-3}$: 26 steps
(6 rejected) reach $3.2$--$3.3 \times 10^{-4}$ on every mesh, in 25~s at
$r=2$ on the GH200.

Tighter error tolerances do not help. Both adaptive variants stall at
$2.3$--$3 \times 10^{-4}$, the spatial error of these meshes, while the step
count doubles with each decade of tolerance. Capped at BDF2, the controller
needs about 30\% more steps for the same accuracy. The monolithic step,
which adaptive BDF requires, costs about 25\% more than the fixed-point step
at a fixed $\Delta t$ (31 against 25~s at $r=4$, $\Delta t = 0.1$~s).

\begin{figure}[h]
\figfile{pareto_hpmr3d-transient}{0.62}
\caption{HP-MR transient: largest relative power error over the run against
solve time. Each code is scored at its own output times against a spline of
the dense reference. Faint: every run; line: each platform's Pareto front.}
\label{fig:transient}
\end{figure}
